% !PRINTED 46
\weberpage{46}

Among these characters there occurs also the one which is obtained by setting $\beta_1 = 0, \beta_2 = 0 \ldots \beta_\nu = 0$, and which has the value 1 for all elements $A$. This we shall call the \gesperrt{unit character}.

The $n$ characters of $S$ can now themselves in turn be combined into an \gesperrt{Abel}ian group, and indeed into a group isomorphic with $S$.

For by (13) each of the $n$ characters is determined by a system of numbers $\beta_1, \beta_2 \ldots \beta_\nu$ taken with respect to the modulus system $a_1, a_2 \ldots a_\nu$, and this system of numbers $\beta$ is at the same time the system of the indices of a definite element $B$ of $S$, namely of
\weq{14}{B = A_1^{\beta_1}A_2^{\beta_2}\ldots A_\nu^{\beta_\nu},}
so that to each element $B$ of $S$ there corresponds a definite character, which we denote by $\chi_B(A)$ or, omitting $A$, by $\chi_B$. This assignment of the characters $\chi$ to the elements $B$ will change, however, if another basis is taken as the foundation, or if the roots of unity $\varepsilon_1, \varepsilon_2 \ldots \varepsilon_\nu$ are chosen differently.

If $B'$ is a second element of $S$ with the indices $\beta_1', \beta_2' \ldots \beta_\nu'$, we obtain in the same way the character $\chi_{B'}$, and if we now understand by $\chi_{BB'}$ the character which for every $A$ is determined by the formula
\weq{15}{\chi_{BB'}(A) = \varepsilon_1^{\alpha_1(\beta_1+\beta_1')}\varepsilon_2^{\alpha_2(\beta_2+\beta_2')}\ldots\varepsilon_\nu^{\alpha_\nu(\beta_\nu+\beta_\nu')}}
then the characters $\chi$ are thereby joined into a group isomorphic with $S$. The unit element in this \gesperrt{group of characters} is the unit character.

Accordingly the group of the characters can also be represented by a basis, which one obtains by setting
\weq{16}{\chi_1(A) = \varepsilon_1^{\alpha_1},\ \chi_2(A) = \varepsilon_2^{\alpha_2} \ldots \chi_\nu(A) = \varepsilon_\nu^{\alpha_\nu}}
Then every character is contained in the form:
\weq{17}{\chi_B = \chi_1^{\beta_1}\chi_2^{\beta_2}\ldots\chi_\nu^{\beta_\nu},}
and, if $B, B'$ are two elements in $S$, we have:
\weq{18}{\chi_B\chi_{B'} = \chi_{BB'}.}

If $\chi_1, \chi_2$ are any two of the characters and $\chi_1\chi_2$ the one composed from the two, then by (15), for every element $A$,
\weq{19}{\chi_1(A)\chi_2(A) = \chi_1\chi_2(A).}
